% Ch.4 Data (~6-8 pp). Drafted Day 1 (full draft Day 2).
% Sourced directly from data.xlsx (sheets: yields, fx, inflation, gdp,
% curve_translation, overview) and the construction in data_preparation.R.
\chapter{Data}
\label{ch:data}

This chapter describes the data used to construct the cycle factors and to run
the predictive regressions of \Cref{ch:results}. All series are assembled in a
single workbook (\texttt{data.xlsx}) and processed into a monthly
country--maturity panel; the construction follows the steps in
\Cref{ch:framework} and is implemented in \texttt{data\_preparation.R}. We
describe the cross-sectional universe, the four raw inputs (yields, inflation,
GDP, and exchange rates) with their sources and units, the sample coverage, and
the cleaning and alignment choices, before turning to summary statistics.

% ===========================================================================
\section{Universe and Sources}
\label{sec:data-sources}

The cross-section is the G10 government-bond universe, comprising eleven
countries: Belgium (BE), Canada (CA), Switzerland (CH), Germany (DE), France
(FR), the United Kingdom (GB), Italy (IT), Japan (JP), the Netherlands (NL),
Sweden (SE), and the United States (US). Five of these share the euro (BE, DE,
FR, IT, NL); the remaining currencies are the Canadian dollar (CAD), Swiss franc
(CHF), pound sterling (GBP), Japanese yen (JPY), Swedish krona (SEK), and US
dollar (USD), the last serving as the numeraire for the international investor.

Four raw inputs enter the analysis, drawn from two data providers:

\begin{description}[leftmargin=1.5em, style=nextline]
  \item[Zero-coupon yields (Bloomberg).] End-of-month yields from Bloomberg's
    fitted zero-coupon government curves (the \texttt{BLC2} family), one curve
    per country, at maturities of $1,2,4,5,9$ and $10$ years. Each country maps
    to a Bloomberg curve identifier (e.g.\ \texttt{I25} for the US, \texttt{I22}
    for the UK, \texttt{I18} for Japan), recorded in the workbook's
    \texttt{curve\_translation} sheet. Yields are continuously compounded and
    quoted in percent per annum.
  \item[Core CPI (LSEG Refinitiv).] Monthly seasonally relevant core consumer
    price indices, one per country (e.g.\ \texttt{aUSCCORF/C} for the US). The
    workbook stores the \emph{index level}; year-on-year inflation
    $\pi_{i,t}=\big(\mathrm{CPI}_{i,t}/\mathrm{CPI}_{i,t-12}-1\big)\times 100$ is
    computed in-sample, and trend inflation $\trendinf_{i,t}$ then follows from
    the EWMA in \eqref{eq:trendinf}.
  \item[Nominal GDP (LSEG Refinitiv).] Annual nominal GDP in local currency
    (e.g.\ \texttt{aUSWDALB5A}), used to form the time-varying global weights
    $w_{i,t}$ of \eqref{eq:gcf}. GDP is converted to a common US-dollar basis
    using year-end exchange rates before weighting.
  \item[Spot exchange rates (Bloomberg).] End-of-month spot rates quoted as US
    dollars per unit of foreign currency (\texttt{EURUSD}, \texttt{CADUSD},
    \texttt{JPYUSD}, \texttt{GBPUSD}, \texttt{SEKUSD}, \texttt{CHFUSD}), with the
    USD rate fixed at unity. These deliver the log currency return
    $s_{i,t+12}-s_{i,t}$ in the dollar-investor excess return \eqref{eq:rx-usd}
    and the FX conversion of GDP.
\end{description}

% ===========================================================================
\section{Sample and Coverage}
\label{sec:data-coverage}

The panel is monthly, sampled at month-end. Although the raw yield workbook
extends back to 1980 and the inflation indices to 1955, the usable span is
governed by the availability of each country's Bloomberg zero-coupon curve. The
United States and Japan enter first, in March and April 1989 respectively; the
remaining nine countries begin in December 1994; all series run through April
2026. This yields $446$ monthly observations for the US, $445$ for Japan, and
$377$ for each of the other countries. \Cref{tab:data-coverage} summarises the
effective coverage.

\begin{table}[t]\centering
  \caption{Effective monthly coverage of the yield panel, by country.}
  \label{tab:data-coverage}
  \small
  \begin{tabular}{llll r}
    \toprule
    Country & Code & Currency & Sample & Months\\
    \midrule
    United States & US & USD & 1989-03 -- 2026-04 & 446\\
    Japan         & JP & JPY & 1989-04 -- 2026-04 & 445\\
    Belgium       & BE & EUR & 1994-12 -- 2026-04 & 377\\
    Canada        & CA & CAD & 1994-12 -- 2026-04 & 377\\
    Switzerland   & CH & CHF & 1994-12 -- 2026-04 & 377\\
    Germany       & DE & EUR & 1994-12 -- 2026-04 & 377\\
    France        & FR & EUR & 1994-12 -- 2026-04 & 377\\
    United Kingdom& GB & GBP & 1994-12 -- 2026-04 & 377\\
    Italy         & IT & EUR & 1994-12 -- 2026-04 & 377\\
    Netherlands   & NL & EUR & 1994-12 -- 2026-04 & 377\\
    Sweden        & SE & SEK & 1994-12 -- 2026-04 & 377\\
    \bottomrule
  \end{tabular}
\end{table}

The panel is therefore unbalanced. This is handled throughout by computing
cross-country objects only over the countries present at each date: the global
weights $w_{i,t}$ in \eqref{eq:gcf} are renormalised each month to sum to one
over the available panel, and the summary correlations below are computed on a
pairwise-complete basis. Spot exchange rates begin in January 1990, so the
US-dollar excess returns of \eqref{eq:rx-usd} for the early US and Japanese
observations (1989) are only available once the corresponding currency series
start; this affects the dollar-investor analysis but not the local-currency
results. \Cref{fig:coverage} visualises the coverage and \Cref{fig:yield-ts} the
resulting 10-year yield series.

\begin{figure}[t]\centering
  \includegraphics[width=0.9\linewidth]{figures/s1_coverage.pdf}
  \caption{Sample coverage by country and maturity. The United States and Japan
    extend back to 1989; the remaining countries begin in late 1994.}
  \label{fig:coverage}
\end{figure}

\begin{figure}[t]\centering
  \includegraphics[width=0.9\linewidth]{figures/s1_yield_ts.pdf}
  \caption{Ten-year zero-coupon yields across the G10, 1989--2026. The secular
    decline in nominal yields and the post-2021 reversal are common across
    countries.}
  \label{fig:yield-ts}
\end{figure}

% ===========================================================================
\section{Cleaning and Variable Construction}
\label{sec:data-cleaning}

The raw inputs are aligned to a common monthly grid keyed on year--month, and
the panel is reshaped to long form (country $\times$ maturity $\times$ month).
Four construction steps deserve note, each tied to \Cref{ch:framework}:

\begin{enumerate}[leftmargin=1.6em]
  \item \textbf{Trend inflation.} Year-on-year core inflation is computed from
    the CPI index and smoothed by the backward-looking EWMA \eqref{eq:trendinf}
    with window $M=120$ months and decay $v=0.987$. Because the CPI history
    predates the yield sample by decades, trend inflation is fully populated
    from the first yield observation, with no warm-up loss inside the estimation
    window.
  \item \textbf{Cycles.} For each country--maturity pair, yields are projected on
    trend inflation by OLS \eqref{eq:decomp}; the residual is the cycle
    $\cyc^{(n)}_{i,t}$. The one-year cycle and the average cycle over
    $\{2,4,5,9,10\}$ years \eqref{eq:cbar} are the inputs to the cycle factor.
  \item \textbf{Excess returns.} One-year holding-period excess returns
    \eqref{eq:rx} are formed for $n\in\{2,5,10\}$ (the maturities for which the
    adjacent shorter yield $y^{(n-1)}$ is observed), duration-standardised by
    $D^{(n)}=n$, and averaged \eqref{eq:rxbar}. Dollar-investor returns add the
    realised currency return and net the US short rate \eqref{eq:rx-usd}.
  \item \textbf{GDP weights.} Annual nominal GDP is converted to US dollars at
    the year-end exchange rate and divided by the cross-sectional total to give
    $w_{i,t}$; weights are held constant within a calendar year and renormalised
    over the available panel.
\end{enumerate}

Missing values arising from the unbalanced panel and from the twelve-month leads
and lags in the return construction are propagated as \texttt{NA} and dropped
pairwise in each regression, rather than imputed.

% ===========================================================================
\section{Summary Statistics}
\label{sec:data-summary}

\Cref{tab:dh-t1-summary} reports the mean and standard deviation of yields for
every country--maturity pair, in the style of \citet{dahlquist2013}, Table~1.
Two features stand out and are consistent across countries: average yield curves
slope upward in maturity, and short-maturity yields are more volatile than long
-maturity yields. The average curve shapes are shown in
\Cref{fig:yield-curve-avg}.

\begin{table}[t]
  \caption{Summary statistics of zero-coupon yields by country and maturity}
  \label{tab:dh-t1-summary}
  \centering
  % dh_t1_summary — Summary statistics of zero-coupon yields by country and maturity.
% Numbers transcribed from tables/dh_t1_summary.pdf.
\footnotesize
\setlength{\tabcolsep}{4pt}
\begin{tabular}{l cc cc cc cc cc cc}
  \toprule
  & \multicolumn{2}{c}{$y^{(1)}_t$} & \multicolumn{2}{c}{$y^{(2)}_t$}
  & \multicolumn{2}{c}{$y^{(4)}_t$} & \multicolumn{2}{c}{$y^{(5)}_t$}
  & \multicolumn{2}{c}{$y^{(9)}_t$} & \multicolumn{2}{c}{$y^{(10)}_t$} \\
  \cmidrule(lr){2-3}\cmidrule(lr){4-5}\cmidrule(lr){6-7}\cmidrule(lr){8-9}\cmidrule(lr){10-11}\cmidrule(lr){12-13}
  Country & Mean & SD & Mean & SD & Mean & SD & Mean & SD & Mean & SD & Mean & SD \\
  \midrule
  Belgium        & $1.88$ & $1.86$ & $2.07$ & $1.92$ & $2.46$ & $2.01$ & $2.64$ & $2.04$ & $3.25$ & $2.08$ & $3.37$ & $2.07$ \\
  Germany        & $1.78$ & $1.86$ & $1.89$ & $1.91$ & $2.19$ & $2.02$ & $2.33$ & $2.04$ & $2.83$ & $2.12$ & $2.91$ & $2.08$ \\
  France         & $1.90$ & $1.94$ & $2.04$ & $1.97$ & $2.39$ & $2.01$ & $2.55$ & $2.01$ & $3.12$ & $2.01$ & $3.23$ & $1.99$ \\
  Italy          & $2.61$ & $2.58$ & $2.96$ & $2.65$ & $3.43$ & $2.60$ & $3.64$ & $2.58$ & $4.28$ & $2.46$ & $4.40$ & $2.43$ \\
  Netherlands    & $1.79$ & $1.86$ & $1.95$ & $1.92$ & $2.30$ & $2.01$ & $2.46$ & $2.02$ & $2.99$ & $2.06$ & $3.08$ & $2.04$ \\
  Switzerland    & $0.85$ & $1.36$ & $0.86$ & $1.36$ & $1.10$ & $1.45$ & $1.23$ & $1.48$ & $1.63$ & $1.56$ & $1.71$ & $1.57$ \\
  United Kingdom & $3.07$ & $2.46$ & $3.16$ & $2.43$ & $3.38$ & $2.32$ & $3.47$ & $2.26$ & $3.77$ & $2.10$ & $3.83$ & $2.05$ \\
  Sweden         & $2.26$ & $2.32$ & $2.43$ & $2.40$ & $2.73$ & $2.45$ & $2.85$ & $2.46$ & $3.23$ & $2.49$ & $3.30$ & $2.48$ \\
  Canada         & $2.77$ & $1.90$ & $2.91$ & $1.89$ & $3.21$ & $1.93$ & $3.32$ & $1.93$ & $3.69$ & $1.92$ & $3.75$ & $1.91$ \\
  Japan          & $0.90$ & $1.89$ & $0.97$ & $1.81$ & $1.19$ & $1.84$ & $1.33$ & $1.89$ & $1.76$ & $1.88$ & $1.85$ & $1.84$ \\
  United States  & $3.12$ & $2.38$ & $3.37$ & $2.40$ & $3.73$ & $2.28$ & $3.90$ & $2.24$ & $4.35$ & $2.10$ & $4.44$ & $2.08$ \\
  \bottomrule
\end{tabular}
\tabnotes{This table reports the mean and standard deviation of monthly
  zero-coupon yields $y^{(n)}_t$, in percent per annum, by country and maturity
  $n\in\{1,2,4,5,9,10\}$ years, over the full per-country sample: the United
  States from 1989-03 and Japan from 1989-04; all other countries from 1994-12;
  all through 2026-04 (see \Cref{tab:data-coverage}). Average yields rise with
  maturity in every market; the volatility of yields is similar across
  maturities, with no uniform direction across countries. After
  \citet{dahlquist2013}, Table~1.}

\end{table}

\begin{figure}[t]\centering
  \includegraphics[width=0.8\linewidth]{figures/s1_yield_curve_avg.pdf}
  \caption{Average zero-coupon yield curve by country over the sample.}
  \label{fig:yield-curve-avg}
\end{figure}

\Cref{tab:dh-t1-corr} reports the cross-country correlation matrix of the
ten-year yield, computed pairwise-complete over the full per-country sample.
Yields co-move strongly internationally---correlations are uniformly high and
are tightest within the euro-area bloc (BE, DE, FR, IT, NL)---which is precisely
the empirical regularity that motivates a \emph{global} cycle factor: if the
levels of long yields share a common component across countries, so plausibly
does the predictable component of their returns. This integration hypothesis is
the subject of \Cref{ch:results}.

\begin{table}[t]
  \caption{Cross-country correlation of the ten-year zero-coupon yield}
  \label{tab:dh-t1-corr}
  \centering
  % dh_t1_corr10y — Cross-country correlation of the 10-year zero-coupon yield.
% Numbers transcribed from tables/dh_t1_corr10y.pdf.
\small
\setlength{\tabcolsep}{4pt}
\begin{tabular}{l ccccccccccc}
  \toprule
  Country & BE & DE & FR & IT & NL & CH & GB & SE & CA & JP & US \\
  \midrule
  Belgium        & $1.00$ &        &        &        &        &        &        &        &        &        &        \\
  Germany        & $0.98$ & $1.00$ &        &        &        &        &        &        &        &        &        \\
  France         & $0.99$ & $0.99$ & $1.00$ &        &        &        &        &        &        &        &        \\
  Italy          & $0.89$ & $0.84$ & $0.88$ & $1.00$ &        &        &        &        &        &        &        \\
  Netherlands    & $0.99$ & $1.00$ & $1.00$ & $0.85$ & $1.00$ &        &        &        &        &        &        \\
  Switzerland    & $0.95$ & $0.97$ & $0.95$ & $0.81$ & $0.97$ & $1.00$ &        &        &        &        &        \\
  United Kingdom & $0.94$ & $0.96$ & $0.96$ & $0.84$ & $0.96$ & $0.90$ & $1.00$ &        &        &        &        \\
  Sweden         & $0.95$ & $0.96$ & $0.96$ & $0.90$ & $0.96$ & $0.95$ & $0.95$ & $1.00$ &        &        &        \\
  Canada         & $0.94$ & $0.97$ & $0.95$ & $0.84$ & $0.96$ & $0.95$ & $0.95$ & $0.98$ & $1.00$ &        &        \\
  Japan          & $0.91$ & $0.91$ & $0.92$ & $0.87$ & $0.90$ & $0.85$ & $0.93$ & $0.92$ & $0.89$ & $1.00$ &        \\
  United States  & $0.90$ & $0.94$ & $0.92$ & $0.75$ & $0.93$ & $0.89$ & $0.95$ & $0.91$ & $0.96$ & $0.90$ & $1.00$ \\
  \bottomrule
\end{tabular}
\tabnotes{This table reports pairwise-complete correlations of the monthly
  ten-year zero-coupon yield $y^{(10)}_t$ across countries, computed over the
  full per-country sample (see \Cref{tab:data-coverage}). Column labels are the
  ISO codes of the countries in the first column. Correlations are uniformly
  high and tightest within the euro-area bloc. After \citet{dahlquist2013},
  Table~1.}

\end{table}
